\documentclass[../../main]{subfiles}

\begin{document}

\section{Operações}
\label{sec:matematica:algebra:estruturas-algebricas:operacoes}

Nesta seção introduzimos o conceito de operação como objeto fundamental da álgebra, a partir do qual estruturas algébricas serão construídas.

\subsection{Operações n-árias}

Seja $A$ um conjunto não vazio. Uma \textbf{operação n-ária} em $A$ é uma função
\[
	\sigma : A^n \to A,
\]
onde $n \in \mathbb{N}, n \geq 1$.

O inteiro $n$ é chamado de aridade da operação.

Casos particulares importantes:
\begin{itemize}
	\item $n=1$: operação unária;
	\item $n=2$: operação binária.
\end{itemize}

\subsection{Operações binárias}

Uma operação binária em $A$ é uma função
\[
	\ast : A \times A \to A.
\]

Uma notação comum é escrever $\ast(a,b)$ como $a \ast b$.

\subsection{Observação estrutural}

Uma operação binária em $A$ não é um subconjunto de $A \times A$, mas uma função total definida em todo par ordenado de elementos de $A$.

Assim, a operação já incorpora a propriedade de fechamento, no sentido de que:
\[
	\forall a,b \in A,\quad a \ast b \in A.
\]

\subsection{Leitura algébrica de operações}

Do ponto de vista algébrico, uma operação não é apenas uma função, mas um símbolo primitivo que será interpretado em estruturas.

Isso permite considerar coleções de operações como objetos estruturais independentes do conjunto subjacente.

\subsection{Estruturas algébricas (visão preliminar)}

Uma estrutura algébrica será definida posteriormente como um conjunto munido de uma família de operações
\[
	\Sigma = \{\sigma_i\}_{i \in I},
\]
cada uma com aridade fixa, juntamente com propriedades que restringem seu comportamento.

\end{document}
